% --- 1. INFORMAZIONI GENERALI ---
\section{Informazioni Generali}
\begin{itemize}
	\item \textbf{Azienda Coinvolta:} {Zucchetti}
	\item \textbf{Mezzo di Comunicazione:} {Google Meet}
	\item \textbf{Data:} 2026-04-01
	\item \textbf{Orario di inizio -- fine:} 11:30 -- 11:50
\end{itemize}


\vspace{0.5cm}
\textbf{Referenti Azienda:}
\begin{table}[ht]
	\centering
	\renewcommand{\arraystretch}{1.2}
	\begin{tabularx}{\textwidth}{X l}
		\toprule
		\textbf{Nome e Cognome} & \textbf{Ruolo} \\
		\midrule
		{Gregorio Piccoli} & {Referente Principale} \\
		\bottomrule
	\end{tabularx}
\end{table}

\vspace{0.5cm}
\textbf{Partecipanti Gruppo (Synergo Unit):}
\begin{table}[ht]
	\centering
	\renewcommand{\arraystretch}{1.2}
	\begin{tabularx}{\textwidth}{X c}
		\toprule
		\textbf{Nome e Cognome} & \textbf{Matricola}\\
		\midrule
		Sofia De Blasi & 2111014 \\
		Roberto Mariano Doroftei & 2111031 \\
		Filippo Idiometri & 2035617 \\
		Francesco Marcon & 2110990 \\
		Armando Moda Scarati & 2082864 \\
		Anton Ricardo Rupi & 2054304 \\
		\bottomrule
	\end{tabularx}
\end{table}

% --- 2. ORDINE DEL GIORNO ---
\vspace{0.5cm}
 
La riunione è stata convocata in seguito all'autorizzazione concessa dal
Prof.\ Vardanega con comunicazione del 2026-03-28, con lo scopo di valutare
congiuntamente la possibilità e la sostenibilità di un aumento di una unità
della capienza del capitolato C6, al fine di consentire al gruppo Synergo Unit
di partecipare alla realizzazione del progetto \textit{Second Brain}.
 
Gli argomenti previsti per la riunione sono:
 
\begin{enumerate}
    \item Motivazioni del gruppo a favore del capitolato C6 e a sfavore
          dell'alternativa disponibile (C4).
    \item Punti di forza del gruppo rispetto al capitolato C6 e prospettive
          di sviluppo oltre i requisiti minimi.
    \item Disponibilità temporale dei membri del gruppo e prospettive accademiche.
    \item Disponibilità del proponente nel periodo iniziale del progetto e
          definizione di un piano di lavoro condiviso.
\end{enumerate}

% --- 3. RESOCONTO ---
\section{Resoconto}
 
\subsection{Valutazione Del Capitolato C4 Come Alternativa}
 
\subsubsection*{Discussione:}
Il proponente ha chiesto al gruppo di esporre le motivazioni per cui il capitolato
C4 non viene ritenuto adeguato come alternativa.
 
\subsubsection*{Risposta e accordi con il proponente:}
 
Il gruppo ha illustrato le seguenti motivazioni, emerse dall'analisi condotta
in fase di candidatura:
 
\begin{itemize}
    \item Il dominio applicativo del capitolato C4 risulta vago e scarsamente
          definito rispetto alle aspettative del gruppo, rendendo difficile
          stimare con precisione l'impegno necessario per il raggiungimento
          di un Minimum Viable Product soddisfacente.
    \item Date le tempistiche vincolanti del progetto didattico, la combinazione
          di un dominio poco definito e di requisiti potenzialmente molto ampi
          non offre al gruppo garanzie sufficienti di poter consegnare un
          prodotto di qualità adeguata entro la scadenza prevista.
    \item Il capitolato C6, al contrario, presenta requisiti ben delimitati,
          un dominio tecnico noto al gruppo e margini espliciti per estensioni
          opzionali, caratteristiche che lo rendono più adatto al contesto
          e alle risorse disponibili.
\end{itemize}
 
\subsection{Punti Di Forza Del Gruppo Rispetto Al Capitolato C6}
 
\subsubsection*{Discussione:}
Il proponente ha invitato il gruppo a illustrare le proprie motivazioni e
i propri punti di forza in relazione al capitolato C6.
 
\subsubsection*{Risposta e accordi con il proponente:}
 
Il gruppo ha evidenziato i seguenti elementi:
 
\begin{itemize}
    \item La conoscenza pregressa delle tecnologie principali richieste
          dal capitolato (HTML, CSS, Markdown, API OpenAI-compatibili) consente
          di concentrare l'impegno sull'integrazione con il modello linguistico
          piuttosto che sull'apprendimento dello stack tecnologico di base.
    \item Il gruppo ha già svolto due incontri formali con il proponente
          (comunicazione email del 2026-03-16 e riunione su Google Meet del
          2026-03-25), nel corso dei quali ha approfondito i requisiti e chiarito
          gli aspetti tecnici, come documentato nei relativi verbali esterni (\href{https://synergo-unit-unipd.github.io/Documentary-Repo/docs/CC/doc_gestionali/verb_esterni/verb_est_2026_03_16_zucchetti/verb_est_2026_03_16_zucchetti.pdf}{\underline{Verbale esterno 2026-03-16}}, \href{https://synergo-unit-unipd.github.io/Documentary-Repo/docs/CC/doc_gestionali/verb_esterni/verb_est_2026_03_25_zucchetti/verb_est_2026_03_25_zucchetti.pdf}{\underline{Verbale esterno 2026-03-25}}).
    \item Il gruppo ha manifestato interesse a esplorare estensioni oltre i
          requisiti minimi del capitolato, in particolare l'implementazione di
          un chatbot con messaggi precompilati come funzionalità opzionale
          aggiuntiva, nell'ottica di sviluppare un prodotto più completo e
          valorizzante per entrambe le parti, sempre tenendo conto delle tempistiche.
\end{itemize}
 
\subsection{Disponibilità Temporale E Prospettive Accademiche}
 
\subsubsection*{Discussione:}
Il proponente ha chiesto informazioni sulla disponibilità temporale complessiva
dei membri del gruppo, sulla scadenza ultima del progetto e sulle prospettive
accademiche individuali (proseguimento con la laurea magistrale o inserimento
nel mondo del lavoro).
 
\subsubsection*{Risposta e accordi con il proponente:}
 
\begin{itemize}
    \item La scadenza ultima prevista per la consegna del progetto è il
      \textbf{2026-08-28}, come dichiarato nel documento di
      \href{https://synergo-unit-unipd.github.io/Documentary-Repo/docs/CC/doc_gestionali/doc_esterni/preventivo_costi_assunzione_impegni/preventivo_costi_e_assunzione_impegni.pdf}
      {\underline{preventivo costi e assunzione impegni}};
      il limite estremo di completamento previsto dal regolamento didattico
      è fissato al 2026-09-26.
    \item I membri del gruppo hanno esposto le proprie prospettive individuali
          riguardo al percorso post-laurea triennale, indicando le intenzioni
          relative all'eventuale proseguimento con la laurea magistrale o
          all'ingresso nel mondo del lavoro.
\end{itemize}
 
\subsection{Disponibilità Del Proponente Nel Periodo Iniziale E Piano Di Lavoro}
 
\subsubsection*{Discussione:}
Il proponente ha illustrato la propria situazione organizzativa per il periodo
immediatamente successivo all'eventuale aggiudicazione del capitolato,
evidenziando una ridotta disponibilità nel primo mese di progetto.
 
\subsubsection*{Risposta e accordi con il proponente:}
 
\begin{itemize}
    \item Il proponente ha comunicato che nel periodo compreso tra la data
      odierna e la fine di aprile 2026 la propria disponibilità sarà
      necessariamente ridotta, in quanto impegnato in parallelo nella
      supervisione di quattro gruppi di progetto.
    \item A partire da maggio 2026 la disponibilità del proponente aumenterà
          significativamente, consentendo un supporto più continuativo al gruppo.
    \item Il gruppo ha dichiarato la propria disponibilità ad adattare il piano
      di lavoro a questa situazione, gestendo autonomamente le attività del
      primo periodo e riducendo le comunicazioni con il proponente allo
      stretto necessario (es.\ chiarimenti puntuali sui requisiti). Le attività
      pianificate per questo periodo includono: impostazione dell'ambiente di
      sviluppo, analisi dei casi d'uso e stesura dell'analisi dei requisiti.
    \item È stato concordato un approccio di lavoro che prevede un ritmo
          inizialmente più graduale nel mese di aprile, con un'intensificazione
          a partire da maggio 2026, in coincidenza con la maggiore disponibilità
          del proponente.
    \item Il gruppo si impegna a mantenere il proponente aggiornato sull'avanzamento
          delle attività anche durante il periodo a disponibilità ridotta,
          tramite comunicazioni asincrone.
\end{itemize}

% --- 4. DECISIONI FORMALI ---
\section{Decisioni Formali}
Nella seguente tabella sono tracciate le decisioni ufficiali e le azioni
conseguenti approvate congiuntamente.
 \renewcommand{\arraystretch}{1.5}
\begin{longtable}{p{1.2cm} p{7.5cm} p{2cm} c}
    \toprule
    \textbf{Cod.} & \textbf{Decisione e Azione}
        & \textbf{Impatto} & \textbf{Scadenza} \\
    \midrule
    \endhead

    VE-4.1
        & \textbf{Decisione:} Il proponente Zucchetti S.p.A.\ accetta di
          aumentare di una unità la capienza del capitolato C6, consentendo
          al gruppo Synergo Unit di partecipare alla realizzazione del progetto
          \textit{Second Brain}.
          \newline
          \textbf{Azione:} Il proponente comunica formalmente l'accettazione;
          il gruppo avvia le attività preliminari del progetto.
        & Organizzativo
        & Immediata \\[4pt]

    VE-4.2
        & \textbf{Decisione:} Nel periodo aprile 2026, a causa del parallelismo
          con altri gruppi, la disponibilità del proponente sarà ridotta. Il gruppo
          impiegherà questo periodo per le attività che richiedono minore
          interazione esterna: analisi dei casi d'uso e analisi dei requisiti.
          \newline
          \textbf{Azione:} Il gruppo pianifica le attività di aprile in modo
          autonomo, aggiornando il proponente tramite comunicazioni asincrone.
        & Pianificazione
        & Fine aprile 2026 \\[4pt]

    VE-4.3
        & \textbf{Decisione:} A partire da maggio 2026 il proponente aumenterà
          la propria disponibilità; il gruppo intensificherà di conseguenza
          il ritmo di lavoro e la frequenza degli incontri formali.
          \newline
          \textbf{Azione:} Pianificazione del primo incontro formale di
          maggio da concordare tra le parti.
        & Pianificazione
        & Inizio maggio 2026 \\

    \bottomrule
\end{longtable}

% --- 5. PROSSIMO INCONTRO ---
\section{Prossimo Incontro}
 
\begin{itemize}
    \item \textbf{Data e Ora:} Da concordare in base alle disponibilità della proponente.
    \item \textbf{Modalità/Luogo:} Google Meet.
    \item \textbf{Argomenti Previsti:}
    \begin{itemize}
        \item Revisione dell'Analisi dei Requisiti prodotta durante il periodo
              iniziale.
        \item Presentazione dell'architettura di massima e delle scelte
              tecnologiche adottate.
        \item Eventuali chiarimenti emersi durante le attività di aprile.
    \end{itemize}
\end{itemize}
 
\newpage
% --- 6. FIRME ---
\section*{Approvazione e Firme}
Con le seguenti firme, i rappresentanti approvano formalmente il contenuto del presente verbale, i requisiti concordati e le decisioni prese.

\vspace{2.5cm}

\noindent
\begin{minipage}[c]{0.45\textwidth}
	\begin{center}
		\rule{6cm}{0.4pt} \\ 
		\vspace{0.2cm}
		\textbf{Per il Gruppo Synergo Unit} \\
		{Sofia De Blasi} \\
		\textit{Responsabile di Progetto}
	\end{center}
\end{minipage}
\hfill
\begin{minipage}[c]{0.45\textwidth}
	\begin{center}
		\rule{6cm}{0.4pt} \\ 
		\vspace{0.2cm}
		\textbf{Per l'Azienda Zucchetti} \\
		{Gregorio Piccoli} \\
		\textit{Referente Aziendale}
	\end{center}
\end{minipage}

\vspace{2cm}
